\documentclass[10pt,letterpaper]{article}
\usepackage{cornell-notes}

\title{Clause 23.04.03.08.04: Basic String Compare}
\author{}
\date{}
\setDocTitle{Clause 23.04.03.08.04: Basic String Compare}
\setDocSubtitle{C++ 2024}
\setDocOwner{C++ 2024 Cornell Notes}
\setDocDate{\today}
\setCornellCollection{C++ 2024 Cornell Notes}
\setCornellUnitType{Clause}
\setCornellUnitNumber{23.04.03.08.04}
\setCornellUnitTitle{Basic String Compare}
\hypersetup{pdftitle={Clause 23.04.03.08.04: Basic String Compare},pdfsubject={Cornell notes},pdfauthor={}}

\begin{document}
\maketitle
\begin{CornellOverviewBox}{Overview}
\textbf{Classification:} Standard library - Normative\\[2pt]
\textbf{Learning objective:} Understand the rules, terminology, and semantic consequences associated with basic\_string::compare.
\end{CornellOverviewBox}\begin{CornellSummaryBox}{Summary}
{\sffamily\bfseries\color{Accent} Summary}\par\vspace{3pt}Section 23.4.3.8.4 focuses on basic\_string::compare. Apply the specified interface together with its mandates, effects, result conditions, exception behavior, and complexity constraints. When applying it, preserve the distinction between mandatory requirements, recommendations, permissions, and behavior deliberately left undefined, unspecified, or implementation-defined.
\end{CornellSummaryBox}

\end{document}
